\NormalFont\ShortTitle{1 Corinteni}

{\MT Prima scrisoare a lui Pavel către Corinteni


\par }\ChapOne{1}{\PP \VerseOne{1}Pavel, chemat să fie apostol al lui Isus Hristos prin voia lui Dumnezeu, și fratele nostru Sostene,

\VS{2}către adunarea lui Dumnezeu care este în Corint, către cei sfințiți în Hristos Isus, numiți sfinți, și către toți cei ce cheamă numele Domnului nostru Isus Hristos în orice loc, atât al lor, cât și al nostru:

\VS{3}Harul și pacea să vă fie vouă, de la Dumnezeu Tatăl nostru și de la Domnul Isus Hristos.


\par }{\PP \VS{4}Mulțumesc totdeauna Dumnezeului meu pentru harul lui Dumnezeu care v-a fost dat în Hristos Isus,

\VS{5}pentru că în toate v-ați îmbogățit în El, în toată vorbirea și în toată cunoștința,

\VS{6}după cum a fost întărită în voimărturia lui Hristos,

\VS{7}ca să nu rămâneți în urmă cu nici un dar, așteptând descoperirea Domnului nostru Isus Hristos,

\VS{8}care vă va întări și el până la sfârșit, nevinovați în ziua Domnului nostru Isus Hristos.

\VS{9}Dumnezeu este credincios, prin care ați fost chemați la părtășia Fiului său, Isus Hristos, Domnul nostru.


\par }{\PP \VS{10}Vă rog, fraților, în Numele Domnului nostru Isus Hristos, ca toți să vorbiți același lucru și să nu fie între voi nici o deosebire, ci să vă desăvârșiți împreună în același gând și în aceeași judecată.

\VS{11}Căci mi s-a raportat despre voi, frații mei, de către cei care sunt din casa lui Chloe, că sunt certuri între voi.

\VS{12}Vreau să spun că fiecare dintre voi spune: “Eu îl urmez pe Pavel”, “Eu îl urmez pe Apolo”, “Eu îl urmez pe Cefa” și “Eu îl urmez pe Hristos”.

\VS{13}Este Hristos împărțit? Pavel a fost răstignit pentru voi? Sau ați fost botezați în numele lui Pavel?

\VS{14}Mulțumesc lui Dumnezeu că nu am botezat pe niciunul dintre voi, în afară de Crispus și de Gaius,

\VS{15}pentru ca nimeni să nu spună că v-am botezat în numele meu.

\VS{16}(Am botezat și casa lui Ștefanas; în afară de ei, nu știu dacă am mai botezat pe altcineva).

\VS{17}Căci Hristos nu m-a trimis ca să botez, ci ca să propovăduiesc Vestea cea Bună — nu cu înțelepciune de cuvinte, pentru ca crucea lui Hristos să nu fie anulată.

\VS{18}Căci cuvântul crucii este o nebunie pentru cei care mor, dar pentru noi, cei care suntem mântuiți, este puterea lui Dumnezeu.

\VS{19}Căci este scris,

\par }{\Q “Voi distruge înțelepciunea celor înțelepți.

\par }{\QB Voi nimici discernământul celor pricepuți.”


\par }{\PP \VS{20}Unde este înțeleptul? Unde este scribul? Unde este cel care discută în acest veac? Oare nu a făcut Dumnezeu nebună înțelepciunea acestei lumi?

\VS{21}Căci, văzând că, în înțelepciunea lui Dumnezeu, lumea, prin înțelepciunea ei, nu a cunoscut pe Dumnezeu, a fost plăcerea lui Dumnezeu ca, prin nebunia propovăduirii, să mântuiască pe cei ce cred.

\VS{22}Căci iudeii cer semne, grecii caută înțelepciunea,

\VS{23}dar noi Îl predicăm pe Hristos răstignit, care este o piatră de poticnire pentru iudei și o nebunie pentru greci,

\VS{24}dar pentru cei chemați, atât iudei cât și greci, Hristos este puterea lui Dumnezeu și înțelepciunea lui Dumnezeu;

\VS{25}pentru că nebunia lui Dumnezeu este mai înțeleaptă decât oamenii, iar slăbiciunea lui Dumnezeu este mai puternică decât oamenii.


\par }{\PP \VS{26}Căci vedeți, fraților, chemarea voastră: nu mulți sunt înțelepți după trup, nu mulți viteji și nu mulți nobili;

\VS{27}ci Dumnezeu a ales lucrurile nebunești ale lumii, ca să facă de rușine pe cei înțelepți. Dumnezeu a ales lucrurile slabe ale lumii ca să facă de rușine pe cele puternice.

\VS{28}Dumnezeu a ales lucrurile smerite ale lumii, lucrurile disprețuite și lucrurile care nu există, ca să nimicească lucrurile care există,

\VS{29}pentru ca nicio făptură să nu se laude înaintea lui Dumnezeu.

\VS{30}Datorită lui, voi sunteți în Hristos Isus, care ne-a fost făcut înțelepciune de la Dumnezeu, și dreptate, și sfințire, și răscumpărare,

\VS{31}pentru ca, așa cum este scris: “Cine se laudă, să se laude în Domnul”.



\par }\Chap{2}{\PP \VerseOne{1}Când am venit la voi, fraților, n-am venit cu o vorbire sau cu înțelepciune, ca să vă vestesc mărturia lui Dumnezeu.

\VS{2}Căci am hotărât să nu cunosc nimic între voi decât pe Isus Hristos și pe cel răstignit.

\VS{3}Am fost cu voi în slăbiciune, în frică și în mult tremur.

\VS{4}Discursul meu și propovăduirea mea nu erau în cuvinte convingătoare ale înțelepciunii omenești, ci în demonstrația Duhului și a puterii,

\VS{5}pentru ca credința voastră să nu stea în înțelepciunea oamenilor, ci în puterea lui Dumnezeu.


\par }{\PP \VS{6}Noi însă vorbim cu înțelepciune printre cei maturi, dar nu cu înțelepciunea lumii acesteia și nici cu cea a conducătorilor lumii acesteia, care sunt pe cale de dispariție.

\VS{7}Ci noi vorbim înțelepciunea lui Dumnezeu în taină, înțelepciunea care a fost ascunsă, pe care Dumnezeu a rânduit-o mai dinainte de lumi, pentru gloria noastră,

\VS{8}și pe care nu a cunoscut-o niciunul dintre conducătorii acestei lumi. Căci, dacă ar fi știut-o, nu L-ar fi răstignit pe Domnul slavei.

\VS{9}Dar, după cum este scris,

\par }{\Q “Lucruri pe care ochiul nu le-a văzut și urechea nu le-a auzit,

\par }{\QB care nu a intrat în inima omului,

\par }{\QB pe care Dumnezeu le-a pregătit pentru cei care îl iubesc.”


\par }{\PP \VS{10}Dar nouă, Dumnezeu ni le-a descoperit prin Duhul Sfânt. Căci Duhul cercetează toate lucrurile, da, chiar și lucrurile adânci ale lui Dumnezeu.

\VS{11}Căci cine dintre oameni cunoaște lucrurile unui om, dacă nu spiritul omului care este în el? Tot așa, nimeni nu cunoaște lucrurile lui Dumnezeu, în afară de Duhul lui Dumnezeu.

\VS{12}Dar noi nu am primit spiritul lumii, ci Duhul care vine de la Dumnezeu, ca să cunoaștem lucrurile care ne-au fost date de Dumnezeu în mod gratuit.

\VS{13}Și noi vorbim aceste lucruri, nu în cuvintele pe care le învață înțelepciunea omenească, ci în cele pe care le învață Duhul Sfânt, comparând lucrurile spirituale cu cele spirituale.

\VS{14}Or, omul firesc nu primește lucrurile Duhului lui Dumnezeu, pentru că sunt o nebunie pentru el; și nu le poate cunoaște, pentru că ele sunt percepute spiritual.

\VS{15}Dar cel spiritual discerne toate lucrurile, și el însuși nu trebuie să fie judecat de nimeni.

\VS{16}“Căci cine a cunoscut gândul Domnului, ca să îl instruiască?” Dar noi avem mintea lui Hristos.



\par }\Chap{3}{\PP \VerseOne{1}Fraților, n-aș putea să vă vorbesc ca unui duhovnicesc, ci ca unui trupesc, ca unor prunci în Hristos.

\VS{2}V-am hrănit cu lapte, nu cu hrană solidă, pentru că nu erați încă pregătiți. Într-adevăr, nici acum nu sunteți pregătiți,

\VS{3}căci sunteți încă trupești. Căci, în măsura în care între voi există gelozie, certuri și facțiuni, nu sunteți voi carnali și nu umblați voi pe căile oamenilor?

\VS{4}Căci atunci când unul spune: “Eu îl urmez pe Pavel”, iar altul: “Eu îl urmez pe Apolo”, nu sunteți voi carnali?


\par }{\PP \VS{5}Cine este deci Apolo și cine este Pavel, dacă nu niște slujitori prin care ați crezut, și fiecare după cum i-a dat Domnul?

\VS{6}Eu am sădit. Apolo a udat. Dar Dumnezeu a dat creșterea.

\VS{7}Așadar, nici cel care plantează nu este nimic, nici cel care udă, ci Dumnezeu, care dă creșterea.

\VS{8}Or, cel care plantează și cel care udă sunt la fel, dar fiecare își va primi răsplata după munca sa.

\VS{9}Căci noi suntem colaboratori ai lui Dumnezeu. Voi sunteți agricultorii lui Dumnezeu, constructorii lui Dumnezeu.


\par }{\PP \VS{10}După harul lui Dumnezeu, care mi-a fost dat, ca un meșter înțelept am pus o temelie, și altul zidește pe ea. Dar fiecare să fie atent cum zidește pe ea.

\VS{11}Căci nimeni nu poate pune o altă temelie decât cea care a fost pusă, care este Isus Hristos.

\VS{12}Dar dacă cineva construiește pe temelie cu aur, argint, pietre scumpe, lemn, fân, paie sau fân,

\VS{13}lucrarea fiecăruia va fi descoperită. Căci Ziua o va vesti, pentru că se descoperă în foc; și focul însuși va testa ce fel de lucrare este lucrarea fiecăruia.

\VS{14}Dacă va rămâne lucrarea cuiva, pe care a construit-o, va primi o răsplată.

\VS{15}Dacă lucrarea cuiva este arsă, va suferi pierderi, dar el însuși va fi salvat, dar ca prin foc.


\par }{\PP \VS{16}Nu știți voi că sunteți templul lui Dumnezeu și că Duhul lui Dumnezeu locuiește în voi?

\VS{17}Dacă cineva va distruge templul lui Dumnezeu, Dumnezeu îl va distruge, căci templul lui Dumnezeu, care sunteți voi, este sfânt.


\par }{\PP \VS{18}Nimeni să nu se înșele pe sine însuși. Dacă cineva crede că este înțelept între voi în lumea aceasta, să se facă nebun, ca să devină înțelept.

\VS{19}Căci înțelepciunea lumii acesteia este o nebunie înaintea lui Dumnezeu. Căci este scris: “El a prins pe cei înțelepți în viclenia lor”.

\VS{20}Și iarăși: “Domnul cunoaște raționamentele înțelepților, care sunt fără valoare.”

\VS{21}De aceea, nimeni să nu se laude în oameni. Căci toate lucrurile sunt ale voastre,

\VS{22}fie că este vorba de Pavel, de Apolo sau de Cefa, de lume, de viață sau de moarte, de cele prezente sau de cele viitoare. Toate sunt ale voastre,

\VS{23}și voi sunteți ai lui Hristos, iar Hristos este al lui Dumnezeu.



\par }\Chap{4}{\PP \VerseOne{1}Deci, să ne considere omul ca slujitori ai lui Hristos și administratori ai tainelor lui Dumnezeu.

\VS{2}Aici, de altfel, se cere de la administratori să fie găsiți credincioși.

\VS{3}Dar în ceea ce mă privește pe mine, este un lucru foarte neînsemnat să fiu judecat de voi sau de un tribunal omenesc. Da, nici măcar nu mă judec pe mine însumi.

\VS{4}Căci nu știu nimic împotriva mea însumi. Cu toate acestea, nu prin aceasta sunt îndreptățit, ci cel care mă judecă este Domnul.

\VS{5}Așadar, nu judecați nimic înainte de vreme, până când va veni Domnul, care va scoate la lumină atât lucrurile ascunse ale întunericului, cât și va descoperi sfaturile inimilor. Atunci fiecare își va primi lauda de la Dumnezeu.


\par }{\PP \VS{6}Și aceste lucruri, fraților, le-am transferat, în chip figurat, mie însumi și lui Apolo, din pricina voastră, ca să învățați în noi să nu gândiți mai mult decât este scris, ca să nu vă îngâmfați unii împotriva altora.

\VS{7}Căci cine vă deosebește? Și ce aveți voi, pe care nu l-ați primit? Dar dacă l-ați primit, de ce vă lăudați ca și cum nu l-ați fi primit?


\par }{\PP \VS{8}Sunteți deja plini. Ați devenit deja bogați. Ai venit să domnești fără noi. Da, și aș vrea să domnești, pentru ca și noi să domnească împreună cu tine!

\VS{9}Căci mă gândesc că Dumnezeu ne-a arătat pe noi, apostolii, cei din urmă, ca pe niște condamnați la moarte. Căci suntem dați în spectacol lumii, atât îngerilor, cât și oamenilor.

\VS{10}Noi suntem nebuni din cauza lui Hristos, dar voi sunteți înțelepți în Hristos. Noi suntem slabi, dar voi sunteți puternici. Voi aveți onoare, dar noi avem dezonoare.

\VS{11}Chiar și în ceasul acesta flămânzim, ne este sete, suntem goi, suntem bătuți și nu avem o locuință sigură.

\VS{12}Ne chinuim, muncind cu mâinile noastre. Când oamenii ne blestemă, noi binecuvântăm. Fiind persecutați, noi îndurăm.

\VS{13}Fiind defăimați, noi implorăm. Suntem făcuți ca mizeria lumii, murdăria ștearsă de toți, chiar și până acum.


\par }{\PP \VS{14}Nu vă scriu aceste lucruri ca să vă rușinez, ci ca să vă îndemn ca pe niște copii iubiți ai Mei.

\VS{15}Căci, deși aveți zece mii de îndrumători în Hristos, nu aveți mulți părinți. Căci, în Cristos Isus, eu am devenit tatăl vostru prin Buna Vestire.

\VS{16}Vă rog, așadar, să fiți imitatori ai mei.

\VS{17}Din această cauză, v-am trimis pe Timotei, care este copilul meu iubit și credincios în Domnul, care vă va reaminti căile mele care sunt în Cristos, așa cum învăț eu pretutindeni, în orice adunare.

\VS{18}Acum, unii sunt îngâmfați, ca și cum n-aș veni la voi.

\VS{19}Dar voi veni la voi în curând, dacă Domnul vrea. Și voi cunoaște, nu cuvântul celor ce se umflă în pene, ci puterea.

\VS{20}Căci Împărăția lui Dumnezeu nu este în cuvânt, ci în putere.

\VS{21}Ce vreți voi? Să vin la voi cu un toiag, sau cu dragoste și cu un duh de blândețe?



\par }\Chap{5}{\PP \VerseOne{1}Se spune că între voi este o imoralitate sexuală, și o imoralitate sexuală cum nu se numește nici măcar printre neamuri: că unul are pe nevasta tatălui său.

\VS{2}Voi sunteți aroganți și nu v-ați jelit, în schimb, ca cel care a făcut această faptă să fie îndepărtat din mijlocul vostru.

\VS{3}Căci eu, cu siguranță, ca absent în trup, dar prezent în duh, am judecat deja, ca și cum aș fi fost de față, pe cel care a făcut acest lucru.

\VS{4}În numele Domnului nostru Isus Cristos, când vă veți aduna împreună cu spiritul meu cu puterea Domnului nostru Isus Cristos,

\VS{5}trebuie să predați pe unul ca acesta Satanei pentru distrugerea cărnii, pentru ca spiritul să fie salvat în ziua Domnului Isus.


\par }{\PP \VS{6}Lăudăroșenia ta nu este bună. Nu știi că puțină drojdie leagă toată făina?

\VS{7}Curățați drojdia veche, ca să fiți un aluat nou, ca și voi, care sunteți azime. Căci, într-adevăr, Cristos, Paștele nostru, a fost sacrificat în locul nostru.

\VS{8}De aceea, să ținem sărbătoarea, nu cu drojdie veche, nici cu drojdia răutății și a răutății, ci cu azimele sincerității și ale adevărului.


\par }{\PP \VS{9}V-am scris în scrisoarea mea să nu vă însoțiți cu păcătoșii sexuali;

\VS{10}dar nicidecum cu păcătoșii sexuali ai lumii acesteia, cu cei lacomi și cu cei ce fac jafuri, cu cei ce fac jafuri și cu cei ce se închină la idoli, pentru că atunci ar trebui să părăsiți lumea.

\VS{11}Dar, așa cum este, v-am scris să nu vă asociați cu nimeni care se numește frate și care este păcătos sexual, sau lacom, sau idolatru, sau calomniator, sau bețiv, sau extorcator. Nici măcar să nu mâncați cu o astfel de persoană.

\VS{12}Căci ce am eu de-a face cu faptul că trebuie să judec și pe cei de afară? Nu-i judecați voi pe cei care sunt înăuntru?

\VS{13}Dar pe cei din afară, Dumnezeu îi judecă. “Îndepărtați-l pe cel rău din mijlocul vostru”.



\par }\Chap{6}{\PP \VerseOne{1}Îndrăznește vreunul din voi, având ceva împotriva aproapelui său, să se ducă la judecată înaintea celor nedrepți și nu înaintea sfinților?

\VS{2}Nu știți că sfinții vor judeca lumea? Și dacă lumea va fi judecată de voi, sunteți voi nevrednici să judecați cele mai mici chestiuni?

\VS{3}Nu știți că noi vom judeca pe îngeri? Cu atât mai mult, lucrurile care țin de această viață?

\VS{4}Dacă deci trebuie să judecați lucruri care țin de viața aceasta, îi puneți să judece pe cei care nu au nici o importanță în adunare?

\VS{5}Spun aceasta ca să vă fac să vă rușinați. Nu există printre voi nici măcar un singur înțelept care să fie în stare să decidă între frații săi?

\VS{6}Dar fratele se judecă cu fratele, și asta în fața unor necredincioși!

\VS{7}Prin urmare, este deja cu totul un defect la voi faptul că aveți procese unii cu alții. De ce să nu fiți mai degrabă nedreptățiți? De ce să nu fiți mai degrabă înșelați?

\VS{8}Nu, ci voi înșivă înșivă faceți nedreptăți și înșelați, și aceasta împotriva fraților voștri.


\par }{\PP \VS{9}Sau nu știți că cei nedrepți nu vor moșteni Împărăția lui Dumnezeu? Nu vă lăsați înșelați. Nici cei imorali din punct de vedere sexual, nici idolatrii, nici adulterii, nici prostituatele, nici homosexualii,

\VS{10}nici hoții, nici lacomii, nici bețivii, nici calomniatorii, nici șantajiștii nu vor moșteni Împărăția lui Dumnezeu.

\VS{11}Unii dintre voi ați fost astfel, dar ați fost spălați. Ați fost sfințiți. Ați fost îndreptățiți în numele Domnului Isus și în Duhul Dumnezeului nostru.


\par }{\PP \VS{12}“Toate lucrurile îmi sunt îngăduite”, dar nu toate sunt potrivite. “Toate lucrurile îmi sunt îngăduite”, dar nu vreau să fiu adus sub puterea a ceva.

\VS{13}“Mâncăruri pentru pântece, și pântecul pentru mâncăruri”, dar Dumnezeu îi va nimici atât pe ea, cât și pe ei. Dar trupul nu este pentru imoralitate sexuală, ci pentru Domnul, și Domnul pentru trup.

\VS{14}Dumnezeu a înviat pe Domnul și ne va învia și pe noi prin puterea Lui.

\VS{15}Nu știți că trupurile voastre sunt mădulare ale lui Hristos? Să iau deci mădularele lui Hristos și să le fac mădulare ale unei prostituate? Să nu se întâmple niciodată așa ceva!

\VS{16}Sau nu știți că cel care este unit cu o prostituată este un singur trup? Căci “Cei doi”, spune el, “vor deveni un singur trup”.

\VS{17}Dar cel care este unit cu Domnul este un singur duh.

\VS{18}Fugiți de imoralitatea sexuală! “Orice păcat pe care îl face un om este în afara trupului”, dar cel care comite imoralitate sexuală păcătuiește împotriva propriului său trup.

\VS{19}Sau nu știți că trupul vostru este un templu al Duhului Sfânt care este în voi, pe care îl aveți de la Dumnezeu? Voi nu sunteți ai voștri,

\VS{20}căci ați fost cumpărați cu un preț. De aceea slăviți pe Dumnezeu în trupul vostru și în duhul vostru, care sunt ale lui Dumnezeu.



\par }\Chap{7}{\PP \VerseOne{1}În ce privește cele despre care mi-ai scris, este bine ca un bărbat să nu se atingă de o femeie.

\VS{2}Dar, din cauza imoralităților sexuale, fiecare bărbat să aibă propria lui soție și fiecare femeie să aibă propriul ei soț.

\VS{3}Bărbatul să dea soției sale afecțiunea care i se cuvine, și tot așa și soția pe soțul ei.

\VS{4}Soția nu are autoritate asupra propriului ei trup, ci soțul. La fel și soțul nu are autoritate asupra propriului său trup, ci soția.

\VS{5}Nu vă lipsiți unul de altul, decât cu consimțământ, pentru o vreme, ca să vă dedicați postului și rugăciunii și să fiți din nou împreună, pentru ca Satana să nu vă ispitească din cauza lipsei voastre de stăpânire de sine.


\par }{\PP \VS{6}Dar aceasta o spun ca o concesie, nu ca o poruncă.

\VS{7}Totuși, aș vrea ca toți oamenii să fie ca mine. Totuși, fiecare om are propriul său dar de la Dumnezeu, unul de acest fel și altul de acel fel.

\VS{8}Dar celor necăsătoriți și văduvelor le spun că este bine pentru ei dacă rămân așa cum sunt eu.

\VS{9}Dar, dacă nu au stăpânire de sine, să se căsătorească. Căci este mai bine să se căsătorească decât să ardă de pasiune.

\VS{10}Celor căsătoriți însă le poruncesc — nu eu, ci Domnul — ca soția să nu-și părăsească soțul

\VS{11}(dar dacă se desparte, să rămână nemăritată sau să se împace cu soțul ei), iar soțul să nu-și părăsească soția.


\par }{\PP \VS{12}Iar celorlalți, nu eu, ci Domnul, le spun: Dacă vreun frate are o nevastă necredincioasă și ea este mulțumită să trăiască cu el, să nu o părăsească.

\VS{13}Femeia care are un soț necredincios, și el este mulțumit să trăiască cu ea, să nu-și părăsească soțul.

\VS{14}Căci soțul necredincios este sfințit în soție, și soția necredincioasă este sfințită în soț. Altfel, copiii voștri ar fi necurați, dar acum sunt sfințiți.

\VS{15}Totuși, dacă cel necredincios se desparte, să fie despărțire. Fratele sau sora nu se află sub robie în astfel de cazuri, ci Dumnezeu ne-a chemat în pace.

\VS{16}Căci de unde știi tu, soție, dacă îți vei salva soțul? Sau de unde știi tu, soț, dacă îți vei salva soția?


\par }{\PP \VS{17}Numai că, după cum a împărțit Domnul fiecăruia, după cum l-a chemat Dumnezeu, așa să umble fiecare. Așa poruncesc eu în toate adunările.


\par }{\PP \VS{18}A fost chemat cineva care a fost tăiat împrejur? Să nu se facă necircumcis. A fost cineva chemat în necircumcizie? Să nu fie circumcis.

\VS{19}Circumcizia nu este nimic, și necircumcizia nu este nimic, dar ceea ce contează este păzirea poruncilor lui Dumnezeu.

\VS{20}Fiecare să rămână în chemarea în care a fost chemat.

\VS{21}Ați fost chemați fiind robi? Nu lăsa ca acest lucru să te deranjeze, dar dacă ai ocazia să devii liber, folosește-o.

\VS{22}Căci cel care a fost chemat în Domnul fiind rob, este omul liber al Domnului. La fel, cel care a fost chemat fiind liber este robul lui Hristos.

\VS{23}Voi ați fost cumpărați cu un preț. Nu vă faceți robi ai oamenilor.

\VS{24}Fraților, fiecare om, în orice condiție a fost chemat, să rămână în acea condiție cu Dumnezeu.


\par }{\PP \VS{25}În ceea ce privește fecioarele, nu am poruncă de la Domnul, ci judec ca unul care a primit îndurare de la Domnul, ca să fie de încredere.

\VS{26}De aceea cred că, din cauza strâmtorării care ne apasă, este bine ca un om să rămână așa cum este.

\VS{27}Ești legat de o soție? Nu căuta să fii eliberat. Ești liber de o soție? Nu căuta o soție.

\VS{28}Dar dacă te căsătorești, nu ai păcătuit. Dacă o fecioară se căsătorește, nu a păcătuit. Cu toate acestea, cei ca ei vor avea opresiune în carne și vreau să vă cruț.

\VS{29}Dar vă spun următorul lucru, fraților: timpul este scurt. De acum încolo, atât cei care au soții pot fi ca și cum nu ar avea niciuna;

\VS{30}și cei care plâng, ca și cum nu ar plânge; și cei care se bucură, ca și cum nu s-ar bucura; și cei care cumpără, ca și cum nu ar avea;

\VS{31}și cei care se folosesc de lume, ca și cum nu s-ar folosi de ea din plin. Căci modul acestei lumi trece.


\par }{\PP \VS{32}Dar eu vreau ca voi să fiți liberi de griji. Cel necăsătorit se îngrijește de lucrurile Domnului, cum să placă Domnului;

\VS{33}dar cel căsătorit se îngrijește de lucrurile lumii, cum să placă soției sale.

\VS{34}Există, de asemenea, o diferență între o soție și o fecioară. Femeia necăsătorită se îngrijește de lucrurile Domnului, ca să fie sfântă atât în trup, cât și în spirit. Dar cea care este căsătorită se îngrijește de lucrurile lumii — cum ar putea să-i placă soțului ei.

\VS{35}Spun aceasta spre folosul vostru, nu ca să vă prind în capcană, ci pentru ceea ce se cuvine și pentru ca voi să vă ocupați de Domnul fără să vă distrageți atenția.


\par }{\PP \VS{36}Dar dacă cineva crede că se poartă necuviincios cu fecioara lui, dacă ea a trecut de vârsta cea mai înaintată, și dacă este nevoie, să facă ce vrea. El nu păcătuiește. Lasă-i să se căsătorească.

\VS{37}Dar cel care stă neclintit în inima sa, fără să aibă nicio urgență, dar care are putere asupra voinței sale și a hotărât în inima sa să-și păstreze fecioara, face bine.

\VS{38}Așadar, atât cel care își dă în căsătorie propria fecioară face bine, cât și cel care nu o dă în căsătorie face mai bine.


\par }{\PP \VS{39}Soția este obligată prin lege atâta timp cât trăiește bărbatul ei; dar dacă soțul a murit, ea este liberă să se mărite cu cine vrea, numai în Domnul.

\VS{40}Dar ea este mai fericită dacă rămâne așa cum este, după părerea mea, și cred că am și eu Duhul lui Dumnezeu.



\par }\Chap{8}{\PP \VerseOne{1}În ceea ce privește lucrurile jertfite idolilor, știm că noi toți avem cunoștință. Cunoașterea umflă, dar dragostea zidește.

\VS{2}Dar dacă cineva crede că știe ceva, nu știe încă așa cum ar trebui să știe.

\VS{3}Dar oricine iubește pe Dumnezeu este cunoscut de el.


\par }{\PP \VS{4}De aceea, cu privire la mâncarea de lucruri jertfite idolilor, știm că nu este niciun idol în lume și că nu este alt Dumnezeu decât unul singur.

\VS{5}Căci, deși există lucruri care se numesc “dumnezei”, fie în ceruri, fie pe pământ — căci sunt mulți “dumnezei” și mulți “domni” —

\VS{6}totuși, pentru noi există un singur Dumnezeu, Tatăl, din care sunt toate lucrurile și noi pentru el; și un singur Domn, Isus Hristos, prin care sunt toate lucrurile și noi trăim prin el.


\par }{\PP \VS{7}Dar această cunoaștere nu este în toți oamenii. Ci unii, având până acum conștiința unui idol, mănâncă ca dintr-un lucru sacrificat unui idol, iar conștiința lor, fiind slabă, este pângărită.

\VS{8}Dar nu mâncarea ne va recomanda lui Dumnezeu. Căci nici dacă nu mâncăm nu suntem mai răi, nici dacă mâncăm nu suntem mai buni.

\VS{9}Dar aveți grijă ca această libertate a voastră să nu devină în niciun caz o piatră de poticnire pentru cei slabi.

\VS{10}Căci, dacă un om vă vede pe voi, care aveți cunoștință, stând în templul unui idol, oare conștiința lui, dacă este slabă, nu va fi încurajată să mănânce lucruri sacrificate idolilor?

\VS{11}Și, prin știința voastră, cel slab piere, fratele pentru care a murit Hristos.

\VS{12}Astfel, păcătuind împotriva fraților și rănindu-le conștiința când este slabă, voi păcătuiți împotriva lui Hristos.

\VS{13}De aceea, dacă mâncarea îl face pe fratele meu să se poticnească, nu voi mai mânca carne în veci, ca să nu-l fac pe fratele meu să se poticnească.



\par }\Chap{9}{\PP \VerseOne{1}Nu sunt eu liber? Nu sunt eu apostol? Nu l-am văzut pe Isus Hristos, Domnul nostru? Nu ești tu lucrarea mea în Domnul?

\VS{2}Dacă pentru alții nu sunt apostol, cel puțin pentru voi sunt, căci voi sunteți pecetea apostolatului meu în Domnul.


\par }{\PP \VS{3}Iată apărarea mea în fața celor ce mă cercetează:

\VS{4}Nu avem noi dreptul să mâncăm și să bem?

\VS{5}Nu avem dreptul să luăm cu noi o soție credincioasă, ca ceilalți apostoli, ca frații Domnului și ca Cefa?

\VS{6}Sau numai eu și Barnaba nu avem dreptul să nu lucrăm?

\VS{7}Ce soldat a servit vreodată pe cheltuiala sa? Cine plantează o vie și nu mănâncă din roadele ei? Sau cine paște o turmă și nu bea din laptele turmei?


\par }{\PP \VS{8}Spun eu oare aceste lucruri după obiceiurile oamenilor? Sau nu spune și Legea același lucru?

\VS{9}Căci este scris în legea lui Moise: “Să nu pui botniță boului în timp ce bate grâul”. Oare pentru boi îi pasă lui Dumnezeu,

\VS{10}sau o spune cu siguranță de dragul nostru? Da, a fost scris de dragul nostru, pentru că cel care ară trebuie să ară în speranță, iar cel care ară în speranță trebuie să aibă parte de speranța lui.

\VS{11}Dacă noi v-am semănat lucruri spirituale, este mare lucru dacă vă secerăm lucruri trupești?

\VS{12}Dacă alții se împărtășesc de acest drept asupra voastră, nu avem noi încă și mai mult?

\par }{\PP Cu toate acestea, nu ne-am folosit de acest drept, ci suportăm totul, ca să nu punem piedici bunei vestiri a lui Hristos.

\VS{13}Nu știți că cei care slujesc în jurul lucrurilor sacre mănâncă din lucrurile templului, iar cei care servesc la altar au partea lor cu altarul?

\VS{14}Tot așa, Domnul a rânduit ca cei care vestesc Buna Vestire să trăiască din Buna Vestire.


\par }{\PP \VS{15}Dar eu n-am folosit nimic din toate acestea și nu scriu aceste lucruri ca să se facă așa în cazul meu, căci mai degrabă aș vrea să mor, decât ca cineva să-mi strice mândria.

\VS{16}Căci, dacă propovăduiesc Vestea cea bună, nu am cu ce să mă laud, pentru că necesitatea îmi este impusă; dar vai de mine dacă nu propovăduiesc Vestea cea bună.

\VS{17}Căci, dacă fac aceasta de bunăvoie, am o răsplată. Dar dacă nu de bunăvoie, am o administrare care mi-a fost încredințată.

\VS{18}Care este, așadar, răsplata mea? Ca, atunci când predic Buna Vestire, să prezint Buna Vestire a lui Cristos fără nicio taxă, pentru a nu abuza de autoritatea mea în Buna Vestire.


\par }{\PP \VS{19}Căci, deși eram liber de toate, m-am supus tuturor, ca să câștig mai mult.

\VS{20}La iudei m-am făcut ca un iudeu, ca să câștig pe iudei; la cei ce sunt sub Lege, ca sub Lege, ca să câștig pe cei ce sunt sub Lege;

\VS{21}la cei fără Lege, ca fără Lege (nefiind fără Lege față de Dumnezeu, ci sub Lege față de Hristos), ca să câștig pe cei fără Lege.

\VS{22}celor slabi m-am făcut ca un slab, ca să câștig pe cei slabi. M-am făcut totul pentru toți, ca să salvez pe unii prin toate mijloacele.

\VS{23}Acum fac aceasta pentru Buna Vestire, ca să fiu părtaș la ea.

\VS{24}Nu știți că cei care aleargă într-o cursă aleargă toți, dar unul singur primește premiul? Aleargă așa, ca să câștigi.

\VS{25}Orice om care se străduiește în jocuri își exercită stăpânirea de sine în toate lucrurile. Or, ei o fac pentru a primi o coroană coruptibilă, dar noi una incoruptibilă.

\VS{26}De aceea, alerg așa, nu fără țintă. Mă lupt așa, nu bătând aerul,

\VS{27}ci îmi bat trupul și îl supun, ca nu cumva, după ce am propovăduit altora, să fiu eu însumi descalificat.



\par }\Chap{10}{\PP \VerseOne{1}Și nu vreau să nu știți, fraților, că părinții noștri au fost cu toții sub nor și au trecut cu toții prin mare,

\VS{2}și că toți au fost botezați în Moise în nor și în mare,

\VS{3}și că toți au mâncat aceeași mâncare duhovnicească,

\VS{4}și că toți au băut aceeași băutură duhovnicească. Căci ei au băut dintr-o stâncă spirituală care i-a urmat, iar stânca era Hristos.

\VS{5}Totuși, cu cei mai mulți dintre ei, Dumnezeu nu a fost mulțumit, pentru că au fost doborâți în pustiu.


\par }{\PP \VS{6}Și aceste lucruri au fost exemplele noastre, ca să nu poftim la lucruri rele, cum au poftit și ei.

\VS{7}Să nu fiți idolatri, cum erau unii dintre ei. După cum este scris: “Poporul ședea să mănânce și să bea și se ridica să se joace”.

\VS{8}Să nu săvârșim imoralități sexuale, cum făceau unii dintre ei, și într-o singură zi au căzut douăzeci și trei de mii de oameni.

\VS{9}Să nu-l punem la încercare pe Hristos, cum au făcut unii dintre ei și au pierit prin șerpi.

\VS{10}Să nu ne cârtiți, cum au cârtit și unii dintre ei și au pierit de distrugător.

\VS{11}Toate aceste lucruri li s-au întâmplat lor ca exemplu și au fost scrise pentru învățătura noastră, asupra cărora au venit sfârșiturile veacurilor.

\VS{12}De aceea, cel care crede că stă în picioare să aibă grijă să nu cadă.


\par }{\PP \VS{13}Nici o ispită nu v-a prins, decât cea obișnuită pentru oameni. Credincios este Dumnezeu, care nu va îngădui să fiți ispitiți peste puterile voastre, ci, odată cu ispita, vă va da și calea de scăpare, ca s-o puteți îndura.


\par }{\PP \VS{14}De aceea, iubiții mei, fugiți de idolatrie.

\VS{15}Vorbesc ca unor oameni înțelepți. Judecați ceea ce spun.

\VS{16}Paharul de binecuvântare pe care îl binecuvântăm, nu este el oare o împărtășire din sângele lui Hristos? Pâinea pe care o frângem, nu este ea o împărtășire din trupul lui Hristos?

\VS{17}Fiindcă există o singură pâine, noi, care suntem mulți, suntem un singur trup, căci toți ne împărtășim din aceeași pâine.

\VS{18}Gândiți-vă la Israel după trup. Cei care mănâncă sacrificiile nu participă oare la altar?


\par }{\PP \VS{19}Atunci ce spun eu? Că un lucru jertfit idolilor este ceva sau că un idol este ceva?

\VS{20}Dar eu spun că lucrurile pe care le sacrifică neamurile le sacrifică demonilor și nu lui Dumnezeu, și nu doresc ca voi să aveți părtășie cu demonii.

\VS{21}Nu puteți bea și paharul Domnului și paharul demonilor. Nu puteți lua parte deopotrivă la masa Domnului și la masa demonilor.

\VS{22}Sau Îl provocăm pe Domnul la gelozie? Suntem noi mai puternici decât el?


\par }{\PP \VS{23}“Toate lucrurile îmi sunt permise, dar nu toate sunt de folos.” “Toate lucrurile îmi sunt permise”, dar nu toate zidesc.

\VS{24}Nimeni să nu-și caute binele propriu, ci fiecare să caute binele aproapelui său.

\VS{25}Tot ce se vinde în măcelărie, să mănânce, fără să pună întrebări din cauza conștiinței,

\VS{26}căci “al Domnului este pământul și plinătatea lui”.

\VS{27}Dar dacă unul dintre cei necredincioși te invită la masă și tu ești înclinat să mergi, mănâncă tot ce ți se pune în față, fără să pui întrebări, din motive de conștiință.

\VS{28}Dar dacă cineva vă spune: “Acest lucru a fost oferit idolilor”, nu-l mâncați, de dragul celui care v-a spus și de dragul conștiinței. Căci “pământul este al Domnului, cu toată plinătatea lui”.

\VS{29}Conștiința, zic, nu a ta, ci a celuilalt. Căci de ce libertatea mea este judecată de o altă conștiință?

\VS{30}Dacă iau parte cu recunoștință, de ce sunt denunțat pentru ceva pentru care mulțumesc?


\par }{\PP \VS{31}De aceea, fie că mâncați, fie că beți, fie că faceți orice altceva, să faceți totul spre slava lui Dumnezeu.

\VS{32}Să nu dați prilej de poticnire, nici iudeilor, nici grecilor, nici adunării lui Dumnezeu;

\VS{33}după cum și eu mulțumesc tuturor în toate, nu caut folosul meu, ci folosul celor mulți, ca să fie mântuiți.



\par }\Chap{11}{\PP \VerseOne{1}Fiți imitatori ai mei, cum și eu sunt al lui Hristos.


\par }{\PP \VS{2}Și acum vă laud, fraților, că vă aduceți aminte de mine în toate lucrurile și că țineți cu tărie tradițiile, așa cum vi le-am dat eu.

\VS{3}Dar vreau să știți că capul oricărui bărbat este Hristos, iar capul femeii este bărbatul, iar capul lui Hristos este Dumnezeu.

\VS{4}Orice om care se roagă sau profețește, având capul acoperit, își necinstește capul.

\VS{5}Dar orice femeie care se roagă sau proorocește cu capul descoperit își dezonorează capul. Căci este unul și același lucru ca și cum ar fi rasă.

\VS{6}Căci, dacă o femeie nu este acoperită, să i se taie și părul. Dar dacă este rușinos pentru o femeie să aibă părul tăiat sau să fie rasă, să fie acoperită.

\VS{7}Căci într-adevăr, bărbatul nu trebuie să aibă capul acoperit, pentru că el este chipul și slava lui Dumnezeu, dar femeia este slava bărbatului.

\VS{8}Căci nu bărbatul vine din femeie, ci femeia vine din bărbat;

\VS{9}căci nu bărbatul a fost creat pentru femeie, ci femeia pentru bărbat.

\VS{10}De aceea, femeia trebuie să aibă autoritate asupra capului ei, din cauza îngerilor.


\par }{\PP \VS{11}Totuși, în Domnul, nici femeia nu este independentă de bărbat, nici bărbatul nu este independent de femeie.

\VS{12}Căci, după cum femeia a ieșit din bărbat, tot așa și bărbatul iese prin femeie; dar toate lucrurile sunt de la Dumnezeu.

\VS{13}Judecați voi înșivă. Se cuvine ca o femeie să se roage lui Dumnezeu fără să fie descoperită?

\VS{14}Oare nu vă învață chiar natura însăși că, dacă un bărbat are părul lung, este o dezonoare pentru el?

\VS{15}Dar dacă o femeie are părul lung, este o glorie pentru ea, căci părul îi este dat ca acoperământ.

\VS{16}Dar dacă vreun bărbat pare a fi certăreț, noi nu avem un astfel de obicei și nici adunările lui Dumnezeu.


\par }{\PP \VS{17}Dar, dându-vă această poruncă, nu vă laud, pentru că vă adunați nu spre bine, ci spre rău.

\VS{18}Pentru că, mai întâi de toate, când vă adunați în adunare, aud că între voi există diviziuni, și în parte cred aceasta.

\VS{19}Căci trebuie să existe și între voi facțiuni, pentru ca cei care sunt aprobați să fie descoperiți printre voi.

\VS{20}De aceea, când vă adunați laolaltă, nu este cina Domnului ceea ce mâncați.

\VS{21}Căci, în timpul mesei voastre, fiecare își ia mai întâi cina lui. Unul este flămând, iar altul este beat.

\VS{22}Ce, n-aveți voi case în care să mâncați și să beți? Sau disprețuiți adunarea lui Dumnezeu și îi faceți de rușine pe cei care nu au destul? Ce să vă spun? Să vă laud? În aceasta nu vă laud.


\par }{\PP \VS{23}Căci am primit de la Domnul ceea ce v-am spus și vouă: că Domnul Isus, în noaptea în care a fost vândut, a luat pâine.

\VS{24}După ce a mulțumit, a frânt-o și a zis:
{\WJ{“Luați, mâncați. Acesta este trupul Meu, care este frânt pentru voi. Faceți aceasta în memoria mea”.
}}

\VS{25}În același fel a luat și paharul după cină, zicând:
{\WJ{“Acest pahar este noul legământ în sângele Meu. Faceți aceasta, de câte ori veți bea, în amintirea mea”.}}

\VS{26}Căci de câte ori mâncați pâinea aceasta și beți paharul acesta, vestiți moartea Domnului până când va veni el.


\par }{\PP \VS{27}De aceea, oricine va mânca pâinea aceasta sau va bea paharul Domnului fără să fie vrednic de Domnul, se va face vinovat de trupul și de sângele Domnului.

\VS{28}Dar omul să se cerceteze pe sine însuși și astfel să mănânce din pâine și să bea din pahar.

\VS{29}Căci cel care mănâncă și bea într-un mod nedemn mănâncă și bea judecata lui însuși, dacă nu discerne trupul Domnului.

\VS{30}Din această cauză mulți dintre voi sunt slabi și bolnavi și nu puțini dorm.

\VS{31}Căci, dacă ne-am discerne pe noi înșine, nu am fi judecați.

\VS{32}Dar când suntem judecați, suntem disciplinați de Domnul, ca să nu fim condamnați împreună cu lumea.

\VS{33}De aceea, frații mei, când vă adunați la masă, așteptați-vă unii pe alții.

\VS{34}Dar dacă cuiva îi este foame, să mănânce acasă, ca nu cumva venirea voastră împreună să fie pentru judecată. Restul, voi pune ordine când voi veni.



\par }\Chap{12}{\PP \VerseOne{1}În ce privește lucrurile spirituale, fraților, nu vreau să fiți neștiutori.

\VS{2}Știți că, atunci când erați păgâni, v-ați lăsat duși de acei idoli muți, oricum ați fi fost conduși.

\VS{3}De aceea vă fac cunoscut că niciun om care vorbește prin Duhul lui Dumnezeu nu spune: “Isus este blestemat”. Nimeni nu poate spune: “Isus este Domnul”, decât prin Duhul Sfânt.


\par }{\PP \VS{4}Și sunt felurite daruri, dar este același Duh.

\VS{5}Există diferite feluri de slujire, și același Domn.

\VS{6}Există diferite feluri de lucrări, dar același Dumnezeu care lucrează totul în toate.

\VS{7}Dar fiecăruia îi este dată manifestarea Duhului, spre folosul tuturor.

\VS{8}Căci unuia i se dă prin Duhul cuvântul înțelepciunii, iar altuia cuvântul cunoștinței, potrivit aceluiași Duh,

\VS{9}altuia credința, prin același Duh, și altuia darurile de vindecare, prin același Duh,

\VS{10}și altuia lucrarea minunilor, și altuia profeția, și altuia discernământul spiritelor, altuia diferite feluri de limbi, și altuia interpretarea limbilor.

\VS{11}Dar unul și același Duh produce toate acestea, împărțind fiecăruia în parte, după cum dorește.


\par }{\PP \VS{12}Căci, după cum trupul este unul și are multe mădulare, și toate mădularele trupului, fiind multe, sunt un singur trup, tot așa și Hristos.

\VS{13}Căci într-un singur Duh am fost botezați toți într-un singur trup, fie iudei, fie greci, fie robi, fie liberi, și toți am fost dați să bem într-un singur Duh.


\par }{\PP \VS{14}Căci trupul nu este un singur mădular, ci mai multe.

\VS{15}Dacă piciorul ar zice: “Pentru că nu sunt mâna, nu fac parte din trup”, nu este deci parte din trup.

\VS{16}Dacă urechea ar spune: “Pentru că nu sunt ochiul, nu fac parte din corp”, nu este deci nu face parte din corp.

\VS{17}Dacă tot corpul ar fi un ochi, unde ar fi auzul? Dacă întregul ar fi auz, unde ar fi mirosul?

\VS{18}Dar acum Dumnezeu a așezat membrele, fiecare dintre ele, în trup, așa cum a dorit.

\VS{19}Dacă toate ar fi un singur mădular, unde ar fi trupul?

\VS{20}Dar acum sunt multe mădulare, dar un singur trup.

\VS{21}Ochiul nu poate spune mâinii: “Nu am nevoie de tine”, sau iarăși capul către picioare: “Nu am nevoie de tine”.

\VS{22}Nu, mult mai degrabă, acele membre ale trupului care par mai slabe sunt necesare.

\VS{23}Acele părți ale trupului pe care le considerăm mai puțin onorabile, pe acelea le acordăm o onoare mai abundentă; și părțile noastre neprezentabile au o modestie mai abundentă,

\VS{24}în timp ce părțile noastre prezentabile nu au o asemenea nevoie. Dar Dumnezeu a alcătuit trupul laolaltă, dând o cinste mai abundentă părții inferioare,

\VS{25}pentru ca în trup să nu existe dezbinare, ci ca membrele să aibă aceeași grijă unele de altele.

\VS{26}Când un membru suferă, toate membrele suferă împreună cu el. Când un membru este onorat, toate membrele se bucură împreună cu el.


\par }{\PP \VS{27}Voi sunteți trupul lui Hristos și sunteți mădularele lui Hristos.

\VS{28}Dumnezeu a așezat pe unii în adunare: întâi apostoli, al doilea profeți, al treilea învățători, apoi făcători de minuni, apoi daruri de vindecare, de ajutor, de guvernare și de diferite feluri de limbi.

\VS{29}Oare toți sunt apostoli? Sunt toți profeți? Sunt toți învățători? Sunt toți făcători de minuni?

\VS{30}Au toți daruri de vindecare? Vorbesc toți în diferite limbi? Sunt toți interpreți?

\VS{31}Dar doriți cu ardoare cele mai bune daruri. Mai mult, vă arăt o cale foarte bună.



\par }\Chap{13}{\PP \VerseOne{1}Dacă vorbesc limbile oamenilor și ale îngerilor, dar nu am dragoste, am ajuns ca o aramă care sună și ca un țambal care sună.

\VS{2}Dacă am darul profeției și cunosc toate tainele și toată știința și dacă am toată credința, încât să mut munții, dar nu am dragoste, nu sunt nimic.

\VS{3}Dacă dau toate bunurile mele pentru a hrăni pe cei săraci și dacă îmi dau trupul pentru a fi ars, dar nu am dragoste, nu-mi folosește la nimic.


\par }{\PP \VS{4}Dragostea este răbdătoare și bună. Dragostea nu invidiază. Iubirea nu se laudă, nu este mândră,

\VS{5}nu se comportă necorespunzător, nu-și caută propria cale, nu se lasă provocată, nu ține cont de rău;

\VS{6}nu se bucură de nedreptate, ci se bucură de adevăr;

\VS{7}suportă totul, crede totul, speră totul și rabdă totul.


\par }{\PP \VS{8}Dragostea nu dă greș niciodată. Dar acolo unde există profeții, ele vor fi desființate. Acolo unde există diverse limbi, ele vor înceta. Acolo unde există cunoaștere, ea va dispărea.

\VS{9}Căci noi cunoaștem în parte și profețim în parte;

\VS{10}dar când va veni ceea ce este complet, atunci ceea ce este parțial va fi desființat.

\VS{11}Când eram copil, vorbeam ca un copil, simțeam ca un copil, gândeam ca un copil. Acum, când am devenit om, am lăsat deoparte lucrurile copilărești.

\VS{12}Căci acum vedem într-o oglindă, slab, dar apoi față în față. Acum cunosc în parte, dar atunci voi cunoaște pe deplin, așa cum și eu am fost cunoscut pe deplin.

\VS{13}Dar acum rămân credința, nădejdea și dragostea — acestea trei. Cea mai mare dintre acestea este dragostea.



\par }\Chap{14}{\PP \VerseOne{1}Urmăriți dragostea și căutați cu râvnă darurile duhovnicești, dar mai ales să prorociți.

\VS{2}Căci cel care vorbește în altă limbă nu vorbește cu oamenii, ci cu Dumnezeu, căci nimeni nu înțelege, ci în Duhul vorbește taine.

\VS{3}Dar cel care profețește vorbește oamenilor pentru zidirea, îndemnul și mângâierea lor.

\VS{4}Cel care vorbește în altă limbă se edifică pe sine însuși, dar cel care profețește edifică adunarea.

\VS{5}Or, eu doresc ca voi toți să vorbiți cu alte limbi, dar și mai mult ca voi să profețiți. Căci mai mare este cel care profețește decât cel care vorbește în alte limbi, dacă nu interpretează, pentru ca adunarea să fie zidită.


\par }{\PP \VS{6}Dar acum, fraților, dacă aș veni la voi vorbind în alte limbi, ce v-aș folosi, dacă nu v-aș vorbi fie prin revelație, fie prin cunoaștere, fie prin profeție, fie prin învățătură?

\VS{7}Chiar și lucrurile lipsite de viață care produc un sunet, fie că este vorba de o cimpoiță sau de o harpă, dacă nu ar da o distincție în sunete, cum s-ar cunoaște ce este cimpoi sau harpă?

\VS{8}Căci, dacă trâmbița ar da un sunet nesigur, cine s-ar pregăti pentru război?

\VS{9}La fel și voi, dacă n-ați rosti prin limbă cuvinte ușor de înțeles, cum s-ar ști ce se vorbește? Căci ați vorbi în aer.

\VS{10}Sunt, poate, atâtea feluri de limbi în lume, și nici una dintre ele nu este lipsită de sens.

\VS{11}Dacă, așadar, n-aș cunoaște sensul limbii, aș fi pentru cel care vorbește un străin, iar cel care vorbește ar fi un străin pentru mine.

\VS{12}Așa și voi, de vreme ce sunteți zeloși pentru darurile spirituale, căutați să abundați pentru zidirea adunării.


\par }{\PP \VS{13}De aceea, cel ce vorbește în altă limbă să se roage ca să poată interpreta.

\VS{14}Căci, dacă mă rog în altă limbă, duhul meu se roagă, dar priceperea mea este nefolositoare.


\par }{\PP \VS{15}Ce ar trebui să fac? Mă voi ruga cu duhul și mă voi ruga și cu mintea. Voi cânta cu duhul și voi cânta și cu mintea.

\VS{16}Altfel, dacă binecuvântezi cu duhul, cum va spune “Amin” la mulțumirea ta cel care ocupă locul celor neînvățați, întrucât nu știe ce spui?

\VS{17}Căci, cu siguranță, tu mulțumești bine, dar celălalt nu este zidit.

\VS{18}Mulțumesc Dumnezeului meu că vorbesc cu alte limbi mai mult decât voi toți.

\VS{19}Totuși, în adunare, prefer să spun cinci cuvinte cu înțelegerea mea, ca să pot instrui și pe alții, decât zece mii de cuvinte în altă limbă.


\par }{\PP \VS{20}Fraților, nu fiți copii în gânduri, dar în răutate fiți copii, ci în gânduri fiți maturi.

\VS{21}În Lege este scris: “Prin oameni de limbi străine și prin buze de străini voi vorbi acestui popor. Nici așa nu mă vor asculta, zice Domnul”.”

\VS{22}De aceea, alte limbi sunt ca un semn, nu pentru cei care cred, ci pentru cei necredincioși; dar profeția este un semn, nu pentru cei necredincioși, ci pentru cei care cred.

\VS{23}Așadar, dacă se adună toată adunarea și toți vorbesc cu alte limbi și intră oameni neînvățați sau necredincioși, nu vor spune ei că sunteți nebuni?

\VS{24}Dar dacă toți propovăduiesc și intră cineva necredincios sau neînvățat, este mustrat de toți și este judecat de toți.

\VS{25}Și astfel se descoperă tainele inimii lui. Astfel, el va cădea cu fața la pământ și se va închina lui Dumnezeu, declarând că Dumnezeu este într-adevăr în mijlocul vostru.


\par }{\PP \VS{26}Atunci, fraților, ce este? Când vă adunați, fiecare dintre voi are un psalm, are o învățătură, are o revelație, are o altă limbă sau are o interpretare. Să faceți totul pentru a vă edifica unii pe alții.

\VS{27}Dacă cineva vorbește într-o altă limbă, să fie doi sau cel mult trei, pe rând, și pe rând, și unul să interpreteze.

\VS{28}Dar, dacă nu este niciun interpret, să tacă în adunare și să vorbească pentru sine și pentru Dumnezeu.

\VS{29}Să vorbească doi sau trei dintre profeți, iar ceilalți să discearnă.

\VS{30}Dar dacă se face o descoperire unui alt om care stă de față, cel dintâi să tacă.

\VS{31}Căci voi toți puteți profeți unul câte unul, ca toți să învețe și toți să fie îndemnați.

\VS{32}Duhurile profeților sunt supuse profeților,

\VS{33}căci Dumnezeu nu este un Dumnezeu al confuziei, ci al păcii, ca în toate adunările sfinților.

\VS{34}Soțiile să facă liniște în adunări, căci nu le-a fost permis să vorbească decât cu supunere, așa cum spune și legea,

\VS{35}dacă doresc să învețe ceva. “Să-și întrebe acasă soții lor, căci este rușinos ca o soție să vorbească în adunare.”

\VS{36}Ce!? Oare de la voi a ieșit cuvântul lui Dumnezeu? Sau a venit numai de la voi?


\par }{\PP \VS{37}Dacă cineva se crede profet sau duhovnic, să recunoască cele ce vă scriu, că sunt porunca Domnului.

\VS{38}Dar dacă cineva este ignorant, să fie ignorant.


\par }{\PP \VS{39}De aceea, fraților, doriți cu tot dinadinsul să prorociți și nu interziceți să vorbiți în alte limbi.

\VS{40}Toate lucrurile să se facă cu decență și în ordine.



\par }\Chap{15}{\PP \VerseOne{1}Și acum, fraților, vă vestesc, fraților, vestea cea bună pe care v-am propovăduit-o, pe care și voi ați primit-o, în care și voi stați în picioare,

\VS{2}prin care și voi sunteți mântuiți, dacă țineți cu tărie cuvântul pe care vi l-am propovăduit, dacă n-ați crezut în zadar.


\par }{\PP \VS{3}Căci v-am dat mai întâi ceea ce am primit și eu: că Hristos a murit pentru păcatele noastre, după Scripturi,

\VS{4}că a fost îngropat, că a înviat a treia zi, după Scripturi,

\VS{5}și că S-a arătat lui Cefa, apoi celor doisprezece.

\VS{6}Apoi s-a arătat deodată la peste cinci sute de frați, dintre care cei mai mulți au rămas până acum, dar și unii au adormit.

\VS{7}Apoi s-a arătat lui Iacov, apoi tuturor apostolilor,

\VS{8}și, în cele din urmă, ca unui copil născut la timpul nepotrivit, mi s-a arătat și mie.

\VS{9}Căci eu sunt cel mai mic dintre apostoli, care nu sunt vrednic să mă numesc apostol, pentru că am prigonit Adunarea lui Dumnezeu.

\VS{10}Dar, prin harul lui Dumnezeu, sunt ceea ce sunt. Harul care mi-a fost dat nu a fost zadarnic, ci am lucrat mai mult decât toți aceștia; dar nu eu, ci harul lui Dumnezeu care era cu mine.

\VS{11}Așadar, fie că sunt eu, fie că sunt ei, așa propovăduim noi și așa ați crezut și voi.


\par }{\PP \VS{12}Și dacă se propovăduiește că Hristos a înviat din morți, cum zic unii dintre voi că nu este înviere a morților?

\VS{13}Dar dacă nu există înviere a morților, nici Hristos nu a înviat.

\VS{14}Dacă Hristos nu a înviat, atunci propovăduirea noastră este zadarnică și credința voastră este și ea zadarnică.

\VS{15}Da, și noi am fost găsiți martori falși ai lui Dumnezeu, pentru că am mărturisit despre Dumnezeu că L-a înviat pe Hristos, pe care nu L-a înviat, dacă este adevărat că morții nu învie.

\VS{16}Căci dacă morții nu au înviat, nici Hristos nu a înviat.

\VS{17}Dacă Hristos nu a înviat, credința voastră este zadarnică; sunteți încă în păcatele voastre.

\VS{18}Atunci și cei care au adormit în Hristos au pierit.

\VS{19}Dacă am nădăjduit în Hristos numai în această viață, suntem dintre toți oamenii cei mai jalnici.


\par }{\PP \VS{20}Dar acum Hristos a înviat din morți. El a devenit cel dintâi rod al celor adormiți.

\VS{21}Căci, de vreme ce moartea a venit prin om, tot prin om a venit și învierea morților.

\VS{22}Căci, după cum în Adam toți mor în Adam, tot așa și în Hristos toți vor fi înviați.

\VS{23}Dar fiecare în ordinea sa: Hristos cel dintâi rod, apoi cei ce sunt ai lui Hristos la venirea Lui.

\VS{24}Apoi vine sfârșitul, când va preda Împărăția lui Dumnezeu Tatăl, când va desființa orice stăpânire și orice autoritate și putere.

\VS{25}Căci trebuie să domnească până când va pune pe toți vrăjmașii Săi sub picioarele Sale.

\VS{26}Ultimul dușman care va fi desființat este moartea.

\VS{27}Căci “El a supus toate lucrurile sub picioarele Sale”. Dar când spune: “Toate lucrurile sunt supuse”, este evident că este exceptat cel care i-a supus toate lucrurile.

\VS{28}După ce toate lucrurile i-au fost supuse, atunci și Fiul însuși va fi supus celui care i-a supus toate lucrurile, pentru ca Dumnezeu să fie totul în toți.


\par }{\PP \VS{29}Altfel, ce vor face cei ce sunt botezați pentru cei morți? Dacă morții nu înviază deloc, atunci de ce sunt botezați pentru cei morți?

\VS{30}De ce stăm și noi în primejdie în fiecare oră?

\VS{31}Afirm că, prin lăudăroșia pe care o am în voi, pe care o am în Hristos Isus, Domnul nostru, mor în fiecare zi.

\VS{32}Dacă m-am luptat cu animalele la Efes pentru scopuri omenești, ce-mi folosește mie? Dacă morții nu învie, atunci “să mâncăm și să bem, căci mâine murim”.

\VS{33}Nu vă lăsați înșelați! “Tovărășiile rele corup bunele moravuri”.

\VS{34}Treziți-vă cu dreptate și nu păcătuiți, căci unii nu au cunoștință de Dumnezeu. Vă spun acest lucru spre rușinea voastră.


\par }{\PP \VS{35}Dar cineva va zice: “Cum înviază morții?” și: “Cu ce fel de trup vin ei?”

\VS{36}Nebunule, ceea ce tu însuți semeni nu se face viu decât dacă moare.

\VS{37}Ceea ce semeni tu însuți, nu semeni trupul care va fi, ci un bob gol, poate de grâu sau de alt fel.

\VS{38}Dar Dumnezeu îi dă un trup, așa cum i-a plăcut lui, și fiecărei semințe un trup propriu.

\VS{39}Nu toate trupurile nu sunt același trup, ci există un trup al oamenilor, un alt trup al animalelor, altul al peștilor și altul al păsărilor.

\VS{40}Există, de asemenea, corpuri cerești și corpuri terestre; dar slava celor cerești diferă de cea a celor terestre.

\VS{41}Există o glorie a soarelui, o altă glorie a lunii și o altă glorie a stelelor; căci o stea diferă de o altă stea în ceea ce privește gloria.


\par }{\PP \VS{42}Tot așa și învierea morților. Trupul este semănat stricăcios, dar înviază nepieritor.

\VS{43}Este semănat în dezonoare; este înviat în glorie. Este semănat în slăbiciune; este înviat în putere.

\VS{44}Este semănat un trup natural; este înviat un trup spiritual. Există un trup natural și există și un trup spiritual.


\par }{\PP \VS{45}Tot așa este scris: “Cel dintâi om, Adam, a devenit un suflet viu.” Ultimul Adam a devenit un duh dătător de viață.

\VS{46}Totuși, nu ceea ce este spiritual este primul, ci ceea ce este natural, apoi ceea ce este spiritual.

\VS{47}Primul om este de pe pământ, făcut din țărână. Al doilea om este Domnul din ceruri.

\VS{48}Cum este cel făcut din țărână, așa sunt și cei care sunt făcuți din țărână; și cum este cel ceresc, așa sunt și cei care sunt cerești.

\VS{49}După cum am purtat chipul celor făcuți din țărână, să purtăm și noi chipul celor cerești.

\VS{50}Acum, fraților, vă spun următoarele: carnea și sângele nu pot moșteni Împărăția lui Dumnezeu; nici ceea ce este pieritor nu moștenește ceea ce este nepieritor.


\par }{\PP \VS{51}Iată, vă spun o taină. Nu toți vom dormi, ci toți vom fi schimbați,

\VS{52}într-o clipă, într-o clipită, într-o clipeală de ochi, la ultima trâmbiță. Căci trâmbița va suna și morții vor fi înviați incoruptibili, iar noi vom fi schimbați.

\VS{53}Căci acest trup pieritor trebuie să devină nepieritor și acest muritor trebuie să se îmbrace în nemurire.

\VS{54}Dar când acest trup perisabil va deveni nepieritor și acest muritor se va îmbrăca în nemurire, atunci se va întâmpla ceea ce este scris: “Moartea este înghițită de biruință”.


\par }{\Q \VS{55}“Moarte, unde este înțepătura ta?

\par }{\QB Hades, unde este victoria ta?”


\par }{\PP \VS{56}Înțepătura morții este păcatul, și puterea păcatului este legea.

\VS{57}Dar mulțumiri fie aduse lui Dumnezeu, care ne dă biruința prin Domnul nostru Isus Hristos.

\VS{58}De aceea, iubiții mei frați, fiți statornici, neclintiți, sporind mereu în lucrarea Domnului, pentru că știți că munca voastră nu este în zadar în Domnul.



\par }\Chap{16}{\PP \VerseOne{1}În ce privește colecta pentru sfinți, așa cum am poruncit adunărilor din Galatia, faceți și voi la fel.

\VS{2}În prima zi a fiecărei săptămâni, fiecare dintre voi să economisească după cum poate prospera, ca să nu se facă colecte când voi veni eu.

\VS{3}Când voi sosi, voi trimite pe oricine veți aproba cu scrisori pentru a duce darul vostru binevoitor la Ierusalim.

\VS{4}Dacă se cuvine să merg și eu, vor merge și ei cu mine.


\par }{\PP \VS{5}Când voi trece prin Macedonia, voi veni la voi, căci trec prin Macedonia.

\VS{6}Dar la voi s-ar putea să rămân la voi, sau chiar să iernez la voi, ca să mă trimiteți la drum, oriunde m-aș duce.

\VS{7}Căci nu vreau să vă văd acum în trecere, ci sper să rămân o vreme cu voi, dacă Domnul va îngădui.

\VS{8}Dar voi rămâne la Efes până la Cincizecime,

\VS{9}pentru că mi s-a deschis o ușă mare și eficientă și sunt mulți adversari.


\par }{\PP \VS{10}Iar dacă vine Timotei, vezi să fie cu voi fără teamă, căci el face lucrarea Domnului, cum fac și eu.

\VS{11}De aceea, nimeni să nu-l disprețuiască. Ci puneți-l înainte, în pace, în călătoria lui, ca să vină la mine, căci îl aștept cu frații.


\par }{\PP \VS{12}În ce privește pe fratele Apolo, l-am îndemnat să vină la voi împreună cu frații, dar nu a dorit deloc să vină acum; dar va veni când va avea ocazia.


\par }{\PP \VS{13}Privește! Rămâneți tari în credință! Fiți curajoși! Fiți tari!

\VS{14}Tot ceea ce faceți să fie făcut cu dragoste.


\par }{\PP \VS{15}Și vă rog, fraților, că știți că casa lui Ștefanas este cea dintâi roadă din Ahaia, și că s-a pus în slujba sfinților,

\VS{16}și voi să vă supuneți lor și tuturor celor ce ajută la lucrare și lucrează.

\VS{17}Mă bucur de venirea lui Ștefanas, a lui Fortunatus și a lui Achaicus; căci ceea ce lipsea din partea voastră, ei au suplinit.

\VS{18}Căci ei au împrospătat spiritul meu și al vostru. De aceea, recunoașteți-i pe cei care sunt astfel.


\par }{\PP \VS{19}Adunările din Asia vă salută. Acuila și Priscila vă salută călduros în Domnul, împreună cu adunarea care este în casa lor.

\VS{20}Toți frații vă salută. Salutați-vă unii pe alții cu o sărutare sfântă.


\par }{\PP \VS{21}Această salutare este făcută de mine, Pavel, cu mâna mea.

\VS{22}Dacă cineva nu iubește pe Domnul Isus Hristos, să fie blestemat. Vino, Doamne!

\VS{23}Harul Domnului Isus Hristos să fie cu voi.

\VS{24}Dragostea mea pentru voi toți în Hristos Isus. Amin.


\par }